% source: 93858040f486e567290db12f5c0581a6b3a1a8c9
\chapter{概率与误差分布}
\label{ch:prob}

\section{什么是分布？}

许多教材会写出这样的表达式：
\[
e \leftarrow \chi,
\]
其中符号 $\chi$ 是一个概率分布。用通俗的话说，这表示值 $e$ 是从某个分布中随机抽样得到的。

例如，掷一枚骰子时，结果是下列集合中的一个：
\[
\{1,2,3,4,5,6\},
\]
每个值以一定的概率出现。这就是一个分布。

形式化地说，分布描述了每种可能结果的概率：
\[
\Pr(X=x).
\]

\section{均匀分布}

最简单的分布是均匀分布。例如，假设该值从下列集合中均匀选取：
\[
\{0,1,2,3\}.
\]
那么每个元素出现的概率都是 $1/4$。

在直方图中，所有柱子的高度都相同。这是最容易理解的随机性概念，但它通常不适合用作 LWE 噪声。

如果噪声在一个很宽的范围上均匀分布，比如从 $-100$ 到 $100$，那么所得误差将过大，解密无法可靠地进行。简而言之，均匀噪声太“野”了。

\section{高斯分布}

许多自然现象近似服从高斯分布，例如人的身高、测量误差以及围绕均值的随机波动。高斯分布以零为中心，呈钟形。标准记法为
\[
N(0,\sigma^2),
\]
表示均值为 $0$、方差为 $\sigma^2$。

参数 $\sigma$ 控制噪声的离散程度。$\sigma$ 较小意味着大多数样本接近零；$\sigma$ 较大意味着样本分布得更分散。

\section{为什么均值通常为零}

在密码学噪声模型中，高斯分布通常以零为中心。原因很简单：噪声不应在某个方向上对系统产生系统性的偏置。如果误差有偏，秘密或解密值就可能泄露信息。

\section{为什么不能在 LWE 中直接使用连续高斯分布}

高斯分布是连续分布。它可以产生如下这样的实数：
\[
0.1352,
\quad -1.2398,
\quad 2.8888.
\]
然而，LWE 定义在有限环或有限域上，通常是模 $q$ 的。误差的各个坐标必须表示为模 $q$ 的整数。

因此，我们使用高斯分布的离散版本。

\section{离散高斯分布}

离散高斯分布（discrete Gaussian）为一个离散的整数集合赋予概率，例如
\[
-3,-2,-1,0,1,2,3.
\]
越靠近零，概率越大；离零越远，概率越小。一种典型的定义是
\[
\Pr(e=x) = \frac{e^{-x^2/(2\sigma^2)}}{Z},
\]
其中 $Z$ 是归一化常数，保证总概率为 $1$。

这类分布之所以重要，是因为它契合基于格的密码的数学结构，并出现在许多安全性证明中。

\section{为什么离散高斯分布如此重要}

离散高斯采样在 Regev 式归约以及格问题的理论分析中处于核心地位。在许多证明中，高斯分布不只是方便而已，而是必不可少的。如果把它换成别的分布，归约论证可能就不成立了。

\section{走向 CBD：Kyber 实际采用的分布}

离散高斯分布在理论上固然优雅，在实践中却颇为棘手：要把它采样好，需要指数运算、浮点比较或拒绝采样循环——这些都很慢，而且都难以做到常数时间，以至于一个草率的采样器会通过其运行时间泄露秘密。出于这一工程上的原因，Kyber 放弃了高斯分布，转而采用\emph{中心二项分布}（centered binomial distribution, CBD）：这是一种仅用比特构造的分布，只需对若干次抛硬币的结果求和即可得到，其钟形足以在安全性上近似高斯分布，同时又能轻而易举地以常数时间采样。

CBD 足够重要——而且足够特定于 ML-KEM——值得单独成章：第~\ref{ch:cbd}~章将完整地展开它，从“分别数出两半中 1 的个数再相减”的做法，一直讲到 FIPS 203 中确切的采样过程及其常数时间性质。这里我们只是把它放在噪声分布演进序列的末端。

\section{$\sigma$ 是什么意思？}

在高斯模型中，参数 $\sigma$ 控制误差的大小。若 $\sigma=0.5$，则大多数误差都非常小；若 $\sigma=100$，则噪声巨大，系统会变得不稳定。在实践中，设计目标是让 $\sigma$ 足够大以隐藏秘密，又足够小以使解密仍能以高概率成功。

\section{尾界}

高斯分布在理论上是无界的，但在实践中，出现极大误差的概率微乎其微。当密码学家谈论一个分布的尾部时，他们指的是取值远离中心的概率。尾部短是有益的，因为这意味着极端误差很罕见，有助于保证正确性。

\section{噪声预算}

每个基于 LWE 和 RLWE 构建的方案都在每个密文中携带一个小的误差项，而正确解密依赖于将该误差保持在某个固定阈值以下——即\emph{噪声预算}（noise budget）。其中的张力是根本性的：噪声选得太小，底层问题就会变得容易，方案因而不安全；噪声选得太大，诚实的解密就会失败。后续章节在为基于 RLWE 的加密和 ML-KEM 确定具体参数时，会再回到这一平衡，那里的设计会把解密失败概率压低到可忽略的水平。管理正确性与安全性之间的余量，是基于格的密码中核心的实际挑战之一。

\section{为什么密码学偏爱高斯噪声}

密码学偏爱高斯噪声，主要有三个原因。

\begin{enumerate}
    \item 它在数学上优雅，便于推理。
    \item 它与基于格的证明和归约自然契合。
    \item 它与最坏情形格困难性之间有很强的理论联系。
\end{enumerate}

简而言之，噪声的存在不只是为了让事情变得随机。它被用来在正确性与安全性之间建立平衡。

\begin{itemize}
    \item 如果噪声太小，秘密就更容易被恢复。
    \item 如果噪声太大，合法用户就无法正确解密。
    \item 各方案中 $q$、$n$、$\eta$、$\sigma$ 等参数的选取，正是为了权衡这些取舍。
\end{itemize}

\section{小结}

本章的思路链条从最简单的随机性一直延伸到 ML-KEM 实际使用的分布：
\[
\text{Uniform} \;\longrightarrow\; \text{Gaussian} \;\longrightarrow\; \text{Discrete Gaussian} \;\longrightarrow\; \text{Centered Binomial Distribution (CBD)}.
\]
均匀噪声易于采样，但对于正确解密而言太“野”；（离散）高斯分布是大多数安全性证明背后在理论上方便的选择；CBD 则是 Kyber 实际采用的折中方案，仅由抛硬币结果之和构造而成。图~\ref{fig:noise-shapes}直接比较了这三种形状：均匀分布的直方图是平坦的，而离散高斯分布和 CBD($\eta=2$) 都已呈现出熟悉的钟形，集中在零附近。

\begin{figure}[H]
\centering
\begin{subfigure}[t]{0.32\textwidth}
\centering
\begin{tikzpicture}
\begin{axis}[
    ybar, bar width=5pt,
    width=\linewidth, height=3.6cm,
    ymin=0, ymax=0.42,
    xtick={-3,-2,-1,0,1,2,3},
    xlabel={$x$}, ylabel={probability},
    label style={font=\scriptsize}, tick label style={font=\scriptsize},
    axis lines=left, enlarge x limits=0.12,
]
\addplot[fill=pqcblue, draw=pqcblue] coordinates {
  (-3,0.143) (-2,0.143) (-1,0.143) (0,0.143) (1,0.143) (2,0.143) (3,0.143)
};
\end{axis}
\end{tikzpicture}
\caption{均匀分布}
\end{subfigure}
\hfill
\begin{subfigure}[t]{0.32\textwidth}
\centering
\begin{tikzpicture}
\begin{axis}[
    ybar, bar width=5pt,
    width=\linewidth, height=3.6cm,
    ymin=0, ymax=0.42,
    xtick={-3,-2,-1,0,1,2,3},
    xlabel={$x$}, ylabel={probability},
    label style={font=\scriptsize}, tick label style={font=\scriptsize},
    axis lines=left, enlarge x limits=0.12,
]
\addplot[fill=pqcteal, draw=pqcteal] coordinates {
  (-3,0.004) (-2,0.054) (-1,0.242) (0,0.399) (1,0.242) (2,0.054) (3,0.004)
};
\end{axis}
\end{tikzpicture}
\caption{离散高斯分布}
\end{subfigure}
\hfill
\begin{subfigure}[t]{0.32\textwidth}
\centering
\begin{tikzpicture}
\begin{axis}[
    ybar, bar width=5pt,
    width=\linewidth, height=3.6cm,
    ymin=0, ymax=0.42,
    xtick={-3,-2,-1,0,1,2,3},
    xlabel={$x$}, ylabel={probability},
    label style={font=\scriptsize}, tick label style={font=\scriptsize},
    axis lines=left, enlarge x limits=0.12,
]
\addplot[fill=pqcorange, draw=pqcorange] coordinates {
  (-2,0.0625) (-1,0.25) (0,0.375) (1,0.25) (2,0.0625)
};
\end{axis}
\end{tikzpicture}
\caption{CBD($\eta=2$)}
\end{subfigure}
\caption{同一范围 $\{-3,\dots,3\}$ 上的三种噪声形状。均匀分布（左，蓝色）是平坦的：每个值出现的可能性相同，这对于可靠解密而言太“野”。离散高斯分布（中，青色）是理论上方便的钟形曲线。$\eta=2$ 的 CBD（右，橙色）仅通过对少量随机比特求和与相减产生（第~\ref{ch:cbd}~章），但其形状已经近似于同一条钟形曲线。}
\label{fig:noise-shapes}
\end{figure}

核心思想是：噪声并不仅仅是一种随机的副产品。它是使 LWE 及相关体制既安全又可解密的主要机制之一。

\printbibliography[heading=chapterrefs]
